\section{Weak compactness and a minimal invariant set}\label{sec:prelim}

We first record the functional-analytic facts needed below, including their
compactness ingredients. The weak topology of a normed space is the topology
generated by its continuous linear functionals. Its continuous dual $X^*$
carries the operator norm, and $X^{**}=(X^*)^*$. Reflexivity means that
the canonical map $J_X:X\to X^{**}$, defined by $(J_Xx)(f)=f(x)$, is onto.
We write
$\overline{\conv} A$ for the norm-closed convex hull of $A$.

\begin{lemma}\label{lem:weak-facts}
Let $X$ be a real normed space.
\begin{enumerate}
\item Every $v\in X$ has a norming functional: there is $f\in X^*$ with
      $\norm{f}\leq 1$ and $f(v)=\norm{v}$. Every norm-closed convex set is
      weakly closed, and the norm is weakly lower semicontinuous.
\item On a finite Cartesian power of $X$, the weak topology is the product
      weak topology. The norm and weak closures of a convex set coincide.
\item If $X$ is reflexive and complete, every norm-closed bounded convex
      subset of $X$ is weakly compact. Arbitrary products of compact
      Hausdorff spaces are compact.
\end{enumerate}
\end{lemma}

The proof, including the compactness details, is given in
Appendix~\ref{app:foundations}.


Throughout, $\N=\{1,2,\ldots\}$. We use the following limit convention.
A proper filter is a family of
subsets closed under finite intersections and supersets, containing the
whole set but not the empty set. An \emph{ultrafilter} is maximal under
inclusion among proper filters. Fix an ultrafilter $\mathcal U$ on $\N$
extending the filter of cofinite sets. It contains no finite set and is
therefore \emph{free}. For a sequence $(a_j)$, the notation
$a=\lim_{j\to\mathcal U}a_j$ means that
$\{j:a_j\in V\}\in\mathcal U$ for every neighborhood $V$ of $a$.
The existence of such $\mathcal U$, and existence and uniqueness of
these limits in compact Hausdorff spaces, are proved in
Appendix~\ref{app:foundations}.

\begin{samepage}
\begin{lemma}\label{lem:reduction}
If Theorem~\ref{thm:main} fails, there are a nonempty separable,
weakly compact convex set $K$ and a nonexpansive map $T:K\to K$ without
a fixed point such that $K$ has no proper nonempty norm-closed convex
$T$-invariant subset. We may take $\diam K=1$ by scalar normalization.
Moreover, the weak topology on $K$ is metrizable. For
\[
          Y=\overline{\operatorname{span}}(K-K),
          \qquad B_Y=\{y\in Y:\norm{y}\leq1\},
\]
the ball $B_Y$ is weakly compact, and $Y$ is separable.
\end{lemma}
\end{samepage}

\begin{proof}
Start with a fixed-point-free map $F:C\to C$ as in the theorem and a point
$a\in C$. Set $C_0=\{a\}$ and
$C_{n+1}=\overline{\conv}(C_n\cup F(C_n))$.
These sets are increasing and separable: continuous images of separable
sets are separable, and rational convex combinations of a countable dense
set are dense in its convex hull. Hence
$C'=\overline{\bigcup_n C_n}$ is separable, bounded, closed, and convex.
Norm continuity of $F$ gives $F(C')\subseteq C'$, and
Lemma~\ref{lem:weak-facts} makes $C'$ weakly compact.

Every decreasing chain of nonempty norm-closed convex invariant subsets
of $C'$ has a nonempty intersection, by weak compactness and the finite
intersection property. The intersection is again closed, convex, and
invariant. Zorn's Lemma, with inclusion reversed, therefore supplies a
minimal such set $K$. Put $T=F|_K$. Its diameter $d$ is positive, since
a singleton invariant set would give a fixed point. From now on express
distances in units of $d$, writing $\norm{\cdot}$ for the original norm
divided by $d$. This scalar normalization preserves the topologies,
nonexpansiveness, and fixed points; it makes $\diam K=1$.

Choose a norm-dense sequence $(z_i)_{i\geq1}$ in $K$. Its pairwise
differences have dense linear span in $Y$, so $Y$ is separable.
Choose a dense sequence $(u_j)$ in the unit sphere of $Y$ and norming
functionals $g_j\in X^*$ for these vectors. They separate points of $Y$:
if $\norm{v}=1$ and $\norm{v-u_j}<1/2$, then $g_j(v)>1/2$.
Thus the map $a\mapsto(g_j(a))_{j\geq1}$ continuously injects weakly
compact $K$ into $\R^{\N}$. The latter is metrizable, for example by
$\sum_j2^{-j}\min\{1,\abs{s_j-t_j}\}$.
The injection is a homeomorphism onto its image, since compact subsets
of a Hausdorff space are closed. Hence the weak topology on $K$ is
metrizable. Finally $B_Y$ is a norm-closed bounded convex subset of $X$,
so is weakly compact by Lemma~\ref{lem:weak-facts}. Its inherited weak
topology equals its intrinsic weak topology: restriction gives one
inclusion of continuous functionals, and the extension argument in
Lemma~\ref{lem:weak-facts} gives the other.
\end{proof}

We argue by contradiction for the remainder of the proof and fix
$K,T,Y$, and $(z_i)$ as in Lemma~\ref{lem:reduction}.
A sequence $(u_j)$ in $K$ is an \emph{approximate fixed point sequence}
if $\norm{Tu_j-u_j}\to0$.
The next lemma is the diameter lemma of Goebel and Karlovitz
\cite{Goebel1975,Karlovitz1976}; we include its proof.
Karlovitz also records independent development of this lemma by
P.~M.~Fitzpatrick \cite[p.~159]{Karlovitz1976}.

\begin{lemma}\label{lem:diametral}
Every approximate fixed point sequence $(u_j)$ satisfies
\[
                     \lim_{j\to\infty}\norm{u_j-z}=1
                     \qquad(z\in K).
\]
\end{lemma}

\begin{proof}
First $K=\overline{\conv}T(K)$: this hull is nonempty and contained in
$K$, and its image under $T$ is contained in $T(K)$ and hence in the hull.
Minimality applies. Consider on $K$ the functions
\[
       s(z)=\sup_{w\in K}\norm{z-w},
       \qquad r(z)=\limsup_j\norm{z-u_j}.
\]
Both are convex and $1$-Lipschitz. For a fixed center, the distance
function is convex and continuous, so its supremum is unchanged on
passing to a norm-closed convex hull. The hull identity and
nonexpansiveness give
\[
 s(Tz)=\sup_{w\in K}\norm{Tz-Tw}\leq s(z),
 \qquad r(Tz)\leq r(z),
\]
where the second inequality also uses $\norm{Tu_j-u_j}\to0$.
Every nonempty sublevel set of either function is norm closed, convex,
and invariant. If the function had two different values, a sublevel
strictly between them would contradict minimality. Both functions
are therefore constant, and $s=1$ because its supremum is $\diam K$.
Write $r=R\leq1$.

Take the weak $\mathcal U$-limit $a\in K$ of $(u_j)$. For every $z\in K$
and $\eta>0$, all sufficiently large $j$ satisfy
$\norm{z-u_j}\leq R+\eta$. The corresponding closed ball is weakly
closed and the cofinite sets belong to $\mathcal U$, so
$\norm{z-a}\leq R+\eta$. Letting $\eta$ decrease to zero and taking
the supremum over $z$ gives $1=s(a)\leq R$, hence $R=1$.
This argument applies to every subsequence of $(u_j)$, since each is
again an approximate fixed point sequence. As all distances are at most
$1$, a failure of the asserted limit would give a subsequence bounded
above by $1-\eta$ for some $\eta>0$, contradicting its limsup of $1$.
\end{proof}

\begin{lemma}\label{lem:compact-detector}
For every nonempty norm-compact set $E\subseteq K$ and every
$\varepsilon>0$, there exists $\delta>0$ such that each $y\in K$ with
$\norm{Ty-y}\leq\delta$ admits a single $f\in X^*$, $\norm{f}\leq1$,
satisfying
\[
                     f(y-z)\geq1-\varepsilon\qquad(z\in E).
\]
\end{lemma}

\begin{proof}
Choose an $\varepsilon/2$-net $a_1,\ldots,a_m$ in $E$ and set
$a=m^{-1}\sum_{i=1}^m a_i\in K$. Lemma~\ref{lem:diametral} implies
that, for some $\delta>0$,
\[
 \norm{Ty-y}\leq\delta
       \quad\Longrightarrow\quad
 \norm{y-a}\geq1-\frac{\varepsilon}{2m}.
\]
Indeed, otherwise choosing a violating point with defect at most $1/j$
for each $j$ would contradict that lemma. Choose a norming functional
$f$ for $y-a$. The numbers $f(y-a_i)$ are all at most $1$, whereas
their sum is $m f(y-a)\geq m-\varepsilon/2$. Each is therefore at least
$1-\varepsilon/2$. For $z\in E$, take $a_i$ with
$\norm{z-a_i}<\varepsilon/2$ and use $\norm{f}\leq1$ to obtain the
claimed bound.
\end{proof}
